\chapter{P2: Konkurensi dan Ketersediaan}
\label{chap:p2}

\section{Rancangan Solusi}

Sistem memisahkan pengambilan data sumber dari pelayanan permintaan pengguna.
Aggregator menjalankan poller BMKG dan PVMBG sebagai dua goroutine independen,
sementara layanan gRPC membaca data kanonis di PostgreSQL. Publisher outbox
berjalan pada goroutine lain. Pemisahan ini terlihat pada
\code{services/aggregator/cmd/main.go}. Dengan demikian, keterlambatan PVMBG
memengaruhi waktu pembaruan data PVMBG, bukan menambahkan waktu tunggu upstream
ke setiap pembacaan BMKG. Permintaan pengguna tetap bergantung pada kapasitas
gateway, autentikasi, Aggregator, dan database yang dipakai bersama.

Setiap poller melakukan pengambilan awal, kemudian mengulangi polling menurut
interval masing-masing. HTTP client memiliki timeout per sumber dan menerima
context untuk pembatalan saat layanan berhenti. Dalam satu poller, satu putaran
diselesaikan sebelum putaran berikutnya diproses; implementasi tidak membuat
goroutine baru tanpa batas pada setiap tick. BMKG mengambil seismic event lalu
peringatan tsunami secara berurutan, sedangkan PVMBG mengambil laporan vulkanik.
Implementasi berada pada \code{services/aggregator/internal/ingest/bmkg.go} dan
\code{services/aggregator/internal/ingest/pvmbg.go}.

\begin{table}[htbp]
    \centering
    \small
    \caption{Parameter konkurensi dan kesegaran data}
    \begin{tabularx}{\textwidth}{@{}P{0.30\textwidth}P{0.18\textwidth}X@{}}
        \toprule
        \textbf{Parameter} & \textbf{Nilai acuan} & \textbf{Makna} \\
        \midrule
        Interval polling kedua sumber & 3.000 ms & Contoh konfigurasi di \code{.env.example}. \\
        Timeout BMKG / PVMBG & 1.000 / 3.500 ms & Batas setiap request HTTP upstream, bukan target latensi API. \\
        Ambang stale & 10.000 ms & Acuan usia keberhasilan pengambilan sumber. \\
        Timeout unary gateway & 2.000 ms & Default deadline panggilan backend. \\
        Permintaan unary aktif & 100 & Default kapasitas semaphore gateway; stream dikecualikan. \\
        Koneksi PostgreSQL & 5--25 & Minimum dan maksimum pool Aggregator. \\
        \bottomrule
    \end{tabularx}
\end{table}

Setiap sumber mempunyai watermark, waktu keberhasilan terakhir, waktu percobaan
terakhir, dan pesan kegagalan. Setelah fetch berhasil, data dinormalisasi dan
disimpan sebelum watermark dimajukan. Pada kegagalan, data sebelumnya tidak
dihapus dan watermark keberhasilan dipertahankan. Pada pembacaan API,
\code{services/aggregator/internal/domain/status.go} menghitung status berikut:
\begin{itemize}
    \item \textbf{FRESH}: keberhasilan terakhir masih dalam ambang waktu dan
    tidak ada error sumber yang tercatat.
    \item \textbf{STALE}: sumber gagal atau terlalu lama tidak berhasil,
    tetapi data kanonis sumber tersebut masih tersedia.
    \item \textbf{UNAVAILABLE}: tidak tersedia data tersimpan untuk digunakan
    ketika sumber belum berhasil atau tidak dapat diperbarui.
\end{itemize}

Respons hazard tetap membawa item yang tersedia beserta metadata
\code{freshness}. Nilai \code{partial=true} diberikan apabila sumber yang
relevan dengan permintaan tidak fresh. Filter sumber membatasi pemeriksaan
kesegaran pada sumber tersebut. Mekanisme ini membedakan kondisi tidak ada
peristiwa dari sumber yang belum dapat memberikan data, sekaligus memungkinkan
klien menampilkan usia data. Pembentukan respons ada pada
\code{services/aggregator/internal/grpc/server.go}; penyimpanan watermark ada
pada \code{services/aggregator/internal/store/hazard_store.go}.

Gateway juga membatasi permintaan unary aktif dengan channel berkapasitas
tetap. Ketika penuh, permintaan ditolak dengan respons overload sesuai kontrak,
umumnya HTTP 429, dan \code{Retry-After}. Deadline unary membatasi waktu tunggu
backend. Pembatas ini bukan pembatas laju per pengguna dan tidak menghitung
stream SSE yang berumur panjang. Kode terkait adalah
\code{services/api-gateway/internal/server/server.go},
\code{services/api-gateway/internal/config/config.go}, dan
\code{services/api-gateway/internal/problem/problem.go}.

\section{Justifikasi dan Alternatif}

Polling independen dipilih karena data kebencanaan dikumpulkan secara berkala
dan dapat dilayani dari penyimpanan kanonis. Pengguna memperoleh respons dari
satu bentuk data meskipun kedua sumber memiliki latensi serta skema berbeda.
Biaya keputusan ini adalah jeda antara perubahan upstream dan kemunculannya
pada API. Karena itu, freshness merupakan bagian dari kontrak respons, bukan
sekadar informasi log.

Alternatif pertama ialah melakukan dua panggilan upstream secara serial pada
setiap request. Pendekatan tersebut sederhana, tetapi latensi sumber terakumulasi
dan gangguan satu sumber langsung memengaruhi request pengguna. Alternatif kedua
ialah fan-out paralel per request dengan deadline terpisah. Pendekatan kedua
mengurangi penjumlahan waktu tunggu, tetapi tetap menambah beban upstream
sebanding dengan jumlah pengguna dan memerlukan fallback tersendiri.

Penyimpanan kanonis dengan polling mengurangi keterikatan tersebut, dengan
konsekuensi kapasitas database menjadi sumber daya bersama. Pool koneksi dan
batas unary gateway menahan pertumbuhan pekerjaan aktif. Penggunaan circuit
breaker atau penjadwal polling terdistribusi dapat menjadi pengembangan
berikutnya, tetapi keduanya belum menjadi mekanisme yang digunakan pada
implementasi ini.

\section{Batasan Solusi}

Isolasi pada tingkat goroutine tidak menjamin isolasi CPU, memori, database,
atau autentikasi. Permintaan BMKG dan PVMBG tetap melewati komponen bersama.
Kegagalan PostgreSQL tidak dapat ditutupi oleh fallback memori karena jalur
baca kanonis masih membutuhkan database. Compose juga belum memberikan bukti
failover beberapa replika Aggregator atau database.

Status kesegaran BMKG terutama mengikuti keberhasilan seismic event. Error
pengambilan peringatan tsunami saat ini diabaikan oleh poller, sehingga
kegagalan pengayaan tsunami belum selalu muncul sebagai status stale.
Selain itu, pembaruan watermark dilakukan sesudah transaksi data selesai,
bukan di dalam transaksi yang sama. Restart pada sela kedua operasi dapat
menyebabkan data diambil kembali; kestabilan identitas hazard membantu menjaga
baris kanonis, tetapi tidak membuktikan tidak ada publication berulang.

Nilai timeout, kapasitas pool, serta semaphore merupakan konfigurasi atau
mekanisme yang ditemukan pada kode. Nilai tersebut tidak membuktikan target
p95, throughput, maupun tingkat availability tertentu tanpa pengukuran.
Tidak ada hasil uji beban gateway terbaru yang dilampirkan untuk snapshot
laporan 9 Oktober 2026.

\section{Proof-of-Concept dan Keterlacakan}

Bukti historis container pada 30 September 2026 disimpan dalam
\code{docs/tex/evidence/mock-container-2026-09-30.json}. Catatan tersebut memuat
10 pemeriksaan berhasil pada lingkup mock. Pengukuran langsung BMKG ketika
empat request PVMBG berjalan paralel tercatat 113,86 ms; pemeriksaan health dan
admin PVMBG tercatat masing-masing 22,29 ms dan 17,80 ms. Angka ini menunjukkan
perilaku pada skenario mock yang dicatat, bukan p95 gateway atau hasil
50 pengguna selama 60 detik.

\begin{table}[htbp]
    \centering
    \small
    \caption{Status pembuktian P2 pada snapshot laporan}
    \begin{tabularx}{\textwidth}{@{}P{0.25\textwidth}X P{0.20\textwidth}@{}}
        \toprule
        \textbf{Skenario} & \textbf{Target atau cakupan} & \textbf{Hasil} \\
        \midrule
        Konkurensi mock & Request BMKG bersamaan dengan empat request PVMBG & Bukti historis tersedia. \\
        Beban API lengkap & 50 VU, 60 detik; error tidak terkendali di bawah 1\%. Skrip juga menetapkan p95 di bawah 2 detik & Belum diukur. \\
        Isolasi melalui gateway & PVMBG ditunda 3 detik; p95 grup BMKG di bawah 300 ms & Belum diukur. \\
        Degradasi sumber & Data lama tetap terbaca dengan freshness yang sesuai & Mekanisme tersedia; hasil terbaru belum dilampirkan. \\
        \bottomrule
    \end{tabularx}
\end{table}

Repositori menyediakan skenario yang dapat diulang pada
\code{tests/load/aggregator.js} dan
\code{tests/load/aggregator-isolation.js}, dengan petunjuk pada
\code{docs/load-testing.md}. Prosedur pembuktian yang direncanakan adalah:
\begin{enumerate}
    \item Jalankan stack yang telah diprovisikan, pastikan akun layanan tersedia
    pada \code{.cache/auth/}, dan buat direktori \code{.cache/load/}.
    Jalankan \code{k6 run tests/load/aggregator.js} dari root repositori.
    Defaultnya adalah pemanasan 10 detik, lalu 50 VU selama 60 detik.
    Campuran request adalah 60\% daftar, 30\% detail, dan 10\% status sumber.
    Hanya fase beban dihitung; kegagalan pengambilan token dicatat terpisah.
    \item Untuk isolasi, tetapkan
    \code{PVMBG_RESPONSE_DELAY_MIN_MS} dan
    \code{PVMBG_RESPONSE_DELAY_MAX_MS} menjadi 3000, lalu buat ulang
    \code{pvmbg-mock} sebagaimana panduan. Jalankan
    \code{k6 run tests/load/aggregator-isolation.js}. Setup memeriksa tiga
    request langsung ke mock dan membatalkan pengujian bila median delay
    kurang dari 2900 ms.
    \item Bandingkan fase baseline 30 detik dengan fase kontensi 60 detik,
    setelah pemanasan 5 detik. Kedua fase menggunakan 10 VU BMKG yang diukur
    dan 10 VU kelompok kedua; kelompok kedua beralih dari BMKG ke PVMBG.
    Simpan p50/p95/p99, error, throughput, ringkasan JSON, dan jendela waktu
    untuk pencocokan trace. Pulihkan konfigurasi delay setelah selesai.
\end{enumerate}

Latensi yang diukur skrip tersebut mencakup Envoy, token exchange, gateway,
RPC, dan pembacaan database. Hasilnya harus dilaporkan sebagai latensi jalur API
lengkap. Keberadaan skrip, threshold, dan instrumen observabilitas merupakan
keterlacakan rancangan pengujian; belum merupakan hasil eksekusi pengujian.

Untuk menilai batas error spesifikasi, HTTP 429 yang merupakan backpressure
terkendali dilaporkan terpisah dari error tidak terkendali. Pada fase beban,
catat jumlah request API $N$, jumlah 429 $N_{429}$, dan kegagalan lain
$N_{\mathrm{err}}$; laporkan $100N_{\mathrm{err}}/N$ serta
$100N_{429}/N$. Denominator dan definisi kegagalan harus dicantumkan agar
angka dapat diperiksa. Threshold error bawaan skrip masih menghitung 429,
sehingga lebih ketat daripada penghitungan penerimaan yang mengecualikannya.
Hasil token endpoint tidak dicampur ke persentil request API.

\begin{table}[htbp]\centering\small
\caption{Format hasil beban yang belum diisi dengan pengukuran}
\begin{tabularx}{\textwidth}{@{}P{.36\textwidth}XX@{}}
\toprule \textbf{Metrik} & \textbf{Beban 50 VU / 60 detik} & \textbf{BMKG saat PVMBG 3 detik}\\\midrule
Jumlah request / throughput & Belum diukur & Belum diukur\\
p50 / p95 / p99 & Belum diukur & Belum diukur\\
Error tidak terkendali / 429 & Belum diukur & Belum diukur\\\bottomrule
\end{tabularx}
\end{table}

Untuk demo outage, catat respons awal kedua jenis hazard, aktifkan
\code{/admin/outage} PVMBG, lalu ulangi request BMKG-only dan volcanic.
Data vulkanik terakhir harus tetap dapat dibaca dengan
\code{stale_since}, atau status unavailable dinyatakan bila data belum ada.
Pertahankan gangguan selama durasi demo yang dipilih, rekam log dan waktu
percobaan, kemudian nonaktifkan outage dan tunggu poller pulih tanpa restart.
Bukti mock historis hanya mencakup 10,64 detik, sehingga tidak membuktikan
skenario outage panjang maupun perilaku fallback API BNPB.
